\chapter{Výběr modelů a vhodnost pro integraci}
\label{ch:selection}

Kapitola~\ref{ch:llm} představila principy fungování velkých jazykových modelů, jejich schopnosti při analýze zdrojového kódu a omezení, která ovlivňují praktické nasazení.
Tato kapitola na tyto poznatky přímo navazuje a zaměřuje se na otázku, jakým způsobem zvolit konkrétní model a formu jeho nasazení v systému Kelvin.

Volba modelu pro školní prostředí nemá pouze technický rozměr.
Kromě schopnosti generovat kvalitní komentáře ke zdrojovému kódu je nutné zohlednit provozní náklady, dostupnost infrastruktury, ochranu osobních údajů studentů a dlouhodobou udržitelnost celého řešení.
Kapitola proto nejprve stanovuje kritéria, podle nichž budou modely a varianty nasazení posuzovány, a následně porovnává dvě základní možnosti nasazení -- cloudové API a on-premise řešení -- z pohledu technické realizace, nákladů, bezpečnosti dat a provozních omezení.

% =============================================================
\section{Kritéria výběru}
\label{sec:selection-criteria}

Na trhu existují desítky modelů s odlišnými vlastnostmi, způsoby nasazení a cenovými modely.
Pro strukturované srovnání je nezbytné nejprve definovat kritéria, která odrážejí specifické požadavky systému Kelvin na analýzu studentských programátorských řešení.
Následující podsekce popisují čtyři klíčová kritéria: kvalitu výstupů, dostupnost modelu, náročnost provozu a schopnost pojmout celé studentské řešení.

\subsection{Kvalita výstupů}
\label{subsec:criteria-quality}

Nejdůležitějším kritériem je schopnost modelu generovat výstupy, které jsou pro učitele i studenty skutečně užitečné.
Kvalita výstupů v kontextu analýzy kódu zahrnuje několik dílčích vlastností:

\begin{itemize}
    \item \textbf{Správnost} -- model musí identifikovat reálné problémy v kódu a nepřidávat falešné výstrahy (\textit{false positives}).
    Falešný pozitivní nález (tedy upozornění na problém, který neexistuje) je pro studenta matoucí a pro učitele zbytečná zátěž.

    \item \textbf{Relevance k zadání} -- model, pokud můžou, by měl posuzovat řešení s ohledem na konkrétní zadání úlohy a ne jen na obecné programátorské principy.
    Například to, zda jde o algoritmus vypisující text na obrazovku v specifickém formátu, není z kódu samotného zjistitelné bez kontextu zadání.

    \item \textbf{Srozumitelnost a didaktická hodnota} -- výstupy musí být psány jazykem, kterému studenti rozumí.
    Model by měl vysvětlit \textit{proč} je daná část kódu problematická a navrhnout konkrétní způsob opravy.

    \item \textbf{Konzistentnost} -- model by měl pro podobné vstupy generovat podobné výstupy.
    Absolutní reprodukovatelnost \textbf{není} u pravděpodobnostních systémů realistická, ale výstupy by se neměly mezi dvěma identickými odevzdáním zásadně lišit.
\end{itemize}

Pro objektivní porovnání kvality modelů existují standardizované benchmarky, což jsou soubory testovacího úloh sloužící k měření výkonu systému za určitých podmínek.
V oblasti generování kódu existuje několik dobře známých benchmarků, které jsou zaměřené na schopnost porozumění a generování kódu.
Nejznámější z nich jsou:

\begin{itemize}
    \item \textbf{HumanEval}~\cite{humaneval} -- sada 164 programátorských úloh v jazyce Python s automatickým ověřením správnosti generovaného kódu.
    Výsledek udává, jak často model na první pokus vygeneruje správné řešení.

    \item \textbf{MBPP} (Mostly Basic Python Problems)~\cite{mbpp} -- podobný benchmark s důrazem na jednoduché úlohy vhodné pro začátečníky.

    \item \textbf{LiveCodeBench}~\cite{livecodebench} -- novější benchmark zahrnující úlohy z programátorských soutěží, které průběžně přibývají,
    čímž snižuje riziko kontaminace trénovaných dat, která nastává u starších benchmarků, které jsou veřejně dostupné již několik let.
    Model by tak mohl mít v trénovaných datech přímo řešení těchto úloh, což by zkreslovalo výsledky.
\end{itemize}

Je důležité mít na paměti, že benchmarky měří především schopnost \textit{generovat} kód, nikoli schopnost \textit{analyzovat} a \textit{komentovat} existující kód.
Proto je nutné, pro hodnocení kvality výstupy v kontextu systému kelvin zvolit jiné metriky, které lépe odrážejí požadavky na analýzu studentských řešení.

\TODO{Doplnit metriky, pravděpodobně založené na hodnocení předešlých odevzdání učiteli a podle toho vyvodit výsledek. Ale nově jsem přidal i přístup, kdy výsledek je následně poslat do dalšího LLM na analýzu. Model ale musí být silnější, aby dokázal posoudit kvalitu výstupu toho prvního LLM.}

\endinput
\subsection{Dostupnost modelu}
\label{subsec:criteria-availability}

Dostupností se v tomto kontextu rozumí několik vzájemně propojených otázek:

\begin{enumerate}
    \item \textbf{Způsob přístupu} -- je model dostupný jako cloudové API nebo jako open-source model ke stažení a lokálnímu nasazení?
    Cloudové API nabízí snadnou integraci bez hardwarových nároků, ale zavádí závislost na externím poskytovateli.
    Open-source modely umožňují lokální nasazení a plnou kontrolu nad zpracováním dat.

    \item \textbf{Licenční podmínky} -- pro školní použití v rámci instituce je nutné ověřit, zda licence modelu takové použití umožňuje.
    Mnoho modelů je dostupných pod striktními licencemi (Apache~2.0, MIT), ale některé mají omezení pro komerční použití nebo vyžadují uzavření samostatné licenční smlouvy.

    \item \textbf{Geografická dostupnost a datová centra} -- jak bylo zmíněno v \secref[Sekci]{subsubsec:gdpr-data-transfer} o ochraně osobních údajů, je důležité, aby poskytovatel modelu měl datová centra v regionech, které jsou v souladu s GDPR, v případě možného úniku dat.

    \item \textbf{Stabilita a podpora} -- API musí být spolehlivé a třetí strana musí poskytovat průběžnou podporu.
    Pro školní použití je důležitá i předvídatelnost cenového modelu a garancí dostupnosti (SLA).
\end{enumerate}

\endinput
\subsection{Náročnost provozu}
\label{subsec:criteria-operational}

Náročnost provozu určuje, jaké zdroje jsou potřeba k dlouhodobému provozování modelu v produkčním prostředí.
Zahrnuje jak finanční náklady, tak nároky na infrastrukturu a údržbu.

U cloudových API jsou náklady typicky určeny cenovými modely a jak bylo popsáno v \secref[Sekci]{subsec:llm-tokens-parameters-context}, jsou založeny na počtu zpracovaných tokenů modelem.
Pro typickou analýzu jednoho studentského odevzdání, zahrnující systémový prompt, zadání a zdrojový kód, je potřeba přibližně 2~000--4~000 vstupních tokenů a 500--2~000 výstupních tokenů při použití one-shot (OS) logiky.
Pro použití chain-of-thought (COT) přístupu, který zahrnuje více kroků a tedy více požadavků na model, se tyto počty mohou zvýšit až na 8~000--16~000 vstupních tokenů a 4~000--8~000 výstupních tokenů.
Celkové náklady na jedno odevzdání tak závisí na zvoleném modelu a jeho cenové politice.

Pro on-premise nasazení jsou náklady primárně jednorázové a zahrnují pořízení hardware, především tedy grafických karet s dostatečnou kapacitou VRAM pro efektivní inference, a náklady na údržbu a aktualizace modelu které jsou průběžně mnohem nižší než u cloudových API, ale vyžadují technickou odbornost pro správu a optimalizaci.

\endinput
\subsection{Schopnost pojmout celé studentské řešení}
\label{subsec:criteria-context-window}

V \secref[Sekci]{subsec:llm-token-limit} byla popsána omezená kapacita kontextového okna LLM, která určuje, kolik tokenů může model zpracovat v jednom požadavku.
Pro analýzu studentských řešení je klíčové, aby model dokázal pojmout celé odevzdání včetně zadání a zdrojového kódu, protože pouze tak může poskytnout relevantní a kontextově správnou zpětnou vazbu.
Zároveň mu musí zůstat dostatek prostoru pro generování kvalitních komentářů.
Je tedy nutné pojmout následující části:

\begin{enumerate}
    \item \textbf{Systémový prompt} -- obsahuje instrukce pro model, které definují jeho roli, požadavky na formát výstupu a další klíčové parametry generování.
    Typicky se jedná o několik strukturovaných vět s jasnými pokyny, které mohou zabrat 200--500 tokenů.

    \item \textbf{Zadání úlohy} -- popis úlohy, kterou student řeší, včetně požadavků na vstupy, výstupy a případné specifické podmínky.
    Zadání může být poměrně rozsáhlé, zejména pokud obsahuje příklady vstupů a výstupů, což může zabrat 500--1~000 tokenů.

    \item \textbf{Zdrojový kód} -- samotné řešení, které student odevzdal.
    Velikost zdrojového kódu se může velmi lišit, ale pro středně velkou úlohu může být v rozsahu 1~000--5~000 tokenů.
\end{enumerate} \TODO{Ověřit konkrétní hodnoty}

Pro standardní studentské odevzdání v systému Kelvin je tedy dostačující kontextové okno přibližně \textbf{8~000 tokenů}.
Aby byl systém schopen zvládnout i rozsáhlejší úlohy (větší projekty, vícesouborová řešení) a aby byl respektován princip dostatečné rezervy, je jako \textbf{minimální požadavek} stanoveno \textbf{16~000 tokenů}.
Moderní modely tuto hranici zpravidla bez problémů splňují, většina aktuálně dostupných modelů nabízí kontextové okno 32~000 až 128~000 tokenů.

\begin{table}[H]
    \centering
    \resizebox{\textwidth}{!}{%
        \begin{tabular}{lp{9cm}l}
            \toprule
            \textbf{Kritérium}              & \textbf{Popis}                                                                 & \textbf{Minimální požadavek} \\
            \midrule
            Kvalita výstupů                 & Věcně správné, srozumitelné, relevantní komentáře                              & Vyšší než průměr \\
            Dostupnost                      & Přístupný přes API nebo ke stažení; vhodná licence pro akademické použití      & Open-source nebo API s DPA \\
            Náročnost provozu               & Přijatelné náklady na token nebo na hardware                                   & Viz \tabref{tab:deployment-comparison} \\
            Kontextové okno                 & Schopnost zpracovat celé odevzdání včetně promptu                              & $\geq$ 16~000 tokenů \\
            \bottomrule
        \end{tabular}
    }
    \caption{Přehled kritérií výběru modelu a jejich minimálních požadavků}
    \label{tab:selection-criteria}
\end{table}

\endinput

% =============================================================
\section{Porovnání možností nasazení}
\label{sec:selection-deployment}

Jazykové modely lze do systému Kelvin integrovat dvěma základními způsoby: jako externí službu poskytovanou prostřednictvím cloudového API, nebo jako lokálně hostovaný model provozovaný přímo na infrastruktuře školy.
Každá z těchto variant přináší odlišné kompromisy mezi jednoduchostí integrace, provozními náklady, ochranou dat a kontrolou nad systémem.
Následující podsekce obě varianty podrobně analyzují z pohledu výše stanovených kritérií.

\subsection{Cloudové API}
\label{subsec:deployment-cloud}

Cloudové API představuje způsob přístupu k jazykovému modelu provozovanému na serverech externího poskytovatele.
Systém Kelvin komunikuje s modelem prostřednictvím HTTP požadavků, tedy odešle vstupní data (prompt, zadání a zdrojový kód) na API endpoint poskytovatele a obdrží vygenerovanou odpověď.
Veškerý výpočet probíhá na straně poskytovatele, takže pro provoz není potřeba žádný specializovaný hardware na školní infrastruktuře.

Mezi hlavní poskytovatele cloudových API pro velké jazykové modely patří:
\begin{itemize}
    \item \textbf{OpenAI} -- modely řady GPT-5.4, GPT-5-mini a další varianty\cite{openai_api}.
    \item \textbf{Anthropic} -- modely Claude Opus~4.6, Claude Sonnet~4.6 a Claude Haiku~4.5\cite{anthropic_api}.
    \item \textbf{Google} -- modely Gemini dostupné prostřednictvím Google AI Studio nebo Vertex AI~\cite{google_vertex}.
    \item \textbf{Mistral AI} -- modely Mistral Large a Mistral Nemo, přičemž některé jsou k dispozici i jako open-source ke stažení~\cite{mistral_api}.
\end{itemize}

\subsubsection*{Hardwarové nároky}
\label{subsubsec:deployment-cloud-hardware}

Z pohledu systému Kelvin jsou hardwarové nároky pro použití cloudového API prakticky nulové.
Na straně Kelvinu je potřeba pouze stabilní internetové připojení a server schopný zpracovávat HTTP požadavky.
Školní infrastruktura, která leží na optické síti CETIN, tuto podmínku bez problémů splňuje.

\subsubsection*{Limity}
\label{subsubsec:deployment-cloud-limits}

Cloudové API zavádějí několik typů omezení, které je nutné při návrhu systému zohlednit:

\begin{enumerate}
    \item \textbf{Limit na počet tokenů v kontextu} -- každý model má maximální délku zpracovatelného vstupu i výstupu.
    U moderních modelů činí tento limit zpravidla 128~000 tokenů nebo více, což je pro potřeby Kelvinu dostatečné.
    Systém však musí zajistit, aby požadavky nepřekročily tento limit, a případnou chybovou odpověď API korektně ošetřit.

    \item \textbf{Limit na počet požadavků} -- poskytovatelé zavádějí takzvané \textit{rate limits}, které omezují počet požadavků v daném časovém období (například 100 požadavků za minutu).
    V kontextu Kelvinu je toto omezení relevantní zejména v období před termínem odevzdání, kdy může počet současných požadavků výrazně vzrůst.
    Systém proto musí implementovat mechanismus řízení zátěže.

    \item \textbf{Finanční náklady} -- cloudové API je zpoplatněno na základě počtu zpracovaných tokenů, což při rozsáhlém nasazení představuje opakující se výdaj.
    Podrobná analýza nákladů pro jednotlivé modely je uvedena níže.

    \item \textbf{Závislost na třetí straně} -- systém Kelvin se stává závislým na dostupnosti a spolehlivosti poskytovatele.
    Výpadek služby, změna API rozhraní nebo ukončení podpory konkrétního modelu mohou narušit funkčnost celého systému.
\end{enumerate}

\subsubsection*{Ochrana dat}
\label{subsubsec:deployment-cloud-data-protection}

Otázka ochrany dat studentů při použití cloudového API byla podrobně rozebrána v~\secref[Sekci]{subsubsec:gdpr-data-transfer}.
V případě cloudového nasazení studentské zdrojové kódy a zadání úloh opouštějí školní infrastrukturu a jsou zpracovávány na serverech poskytovatele.
Tato skutečnost může být rozhodujícím faktorem pro volbu mezi cloudovým API a on-premise nasazením.

\subsubsection*{Náklady}
\label{subsubsec:deployment-cloud-costs}

Náklady cloudových API jsou určeny cenovým modelem \textit{za token}.
Poskytovatelé zpravidla rozlišují vstupní tokeny (text odesílaný modelu) a výstupní tokeny (text generovaný modelem), přičemž výstupní tokeny bývají výrazně dražší.

% ===================================================================
%   Chain-of-thought přístup (3 kroky: Draft analysis, Critique analysis, Review analysis):
%     Průměr vstupních tokenů:  9 053  (min: 7 247, max: 9 687)
%     Průměr výstupních tokenů: 5 149  (min: 3 863, max: 6 404)
%
%   One-shot přístup (1 krok: One-shot analysis):
%     Průměr vstupních tokenů:  2 043  (min:   891, max: 4 508)
%     Průměr výstupních tokenů:   501  (min:   305, max:   590)
% ===================================================================
\newcommand{\chainInToks}{9053}
\newcommand{\chainOutToks}{5149}
\newcommand{\oneshotInToks}{2043}
\newcommand{\oneshotOutToks}{501}

% Výpočet nákladů: #1 = cena vstupu ($/1M tokenů), #2 = cena výstupu ($/1M tokenů)
%                  #3 = počet vstupních tokenů,    #4 = počet výstupních tokenů
\newcommand{\computecost}[4]{%
  \begingroup
  \pgfkeys{/pgf/fpu, /pgf/fpu/output format=fixed}%
  \pgfmathparse{(#3 * (#1 / 1000000)) + (#4 * (#2 / 1000000))}%
  \pgfmathsetmacro{\myresult}{\pgfmathresult}%
  \pgfkeys{/pgf/fpu=false}%
  \pgfkeys{/pgf/number format/.cd, use comma, fixed, precision=4}%
  \pgfmathprintnumber{\myresult}%
  \endgroup
}

% Výpočet nákladů pro Chain-of-Thought
% #1 = cena vstupu ($/1M), #2 = cena výstupu ($/1M)
\newcommand{\computeCostCOT}[2]{%
  \computecost{#1}{#2}{\chainInToks}{\chainOutToks}%
}

% Výpočet nákladů pro One-shot
% #1 = cena vstupu ($/1M), #2 = cena výstupu ($/1M)
\newcommand{\computeCostOS}[2]{%
  \computecost{#1}{#2}{\oneshotInToks}{\oneshotOutToks}%
}

\begin{table}[H]
    \centering
    \resizebox{\textwidth}{!}{%
        \begin{tabular}{lllll}
            \toprule
            & & & \multicolumn{2}{c}{\textbf{Náklady na odevzdání}} \\
            \textbf{Model} & \textbf{Vstup (\$/1M)} & \textbf{Výstup (\$/1M)} & \textbf{chain-of-thought (\$)} & \textbf{one-shot (\$)} \\
            \midrule
                GPT-5.4\cite{gpt-5.4}                               & 2.50  & 15.00  & $\sim\computeCostCOT{2.50}{16.00}$ & $\sim\computeCostOS{2.50}{15.00}$ \\
                GPT-5.4-mini\cite{gpt-5.4-mini}                     & 0.75  & 4.50   & $\sim\computeCostCOT{0.75}{4.50}$  & $\sim\computeCostOS{0.75}{4.50}$ \\
                GPT-5.4-nano\cite{gpt-5.4-nano}                     & 0.20  & 1.25   & $\sim\computeCostCOT{0.20}{1.25}$  & $\sim\computeCostOS{0.20}{1.25}$ \\
            \midrule
                Claude 4.6 Opus\cite{claude-opus-4.6}               & 5.00  & 25.00  & $\sim\computeCostCOT{5.00}{25.00}$ & $\sim\computeCostOS{5.00}{25.00}$ \\
                Claude 4.6 Sonnet\cite{claude-sonnet-4.6}           & 3.00  & 15.00  & $\sim\computeCostCOT{3.00}{15.00}$ & $\sim\computeCostOS{3.00}{15.00}$ \\
                Claude 3.5 Haiku\cite{claude-3.5}                   & 0.80  & 4.00   & $\sim\computeCostCOT{0.80}{4.00}$  & $\sim\computeCostOS{0.80}{4.00}$ \\
                Claude 3.0 Haiku\cite{claude-3}                     & 0.25  & 1.25   & $\sim\computeCostCOT{0.25}{1.25}$  & $\sim\computeCostOS{0.25}{1.25}$ \\
            \midrule
                Gemini 3.1 Pro\cite{geminy-3.1-pro}                 & 2.00  & 12.00  & $\sim\computeCostCOT{2.00}{12.00}$ & $\sim\computeCostOS{2.00}{12.00}$ \\
                Gemini 3.1 Flash-Lite\cite{geminy-3.1-flash-lite}   & 0.25  & 1.50   & $\sim\computeCostCOT{0.25}{1.50}$  & $\sim\computeCostOS{0.25}{1.50}$ \\
                Gemini 3 Flash\cite{geminy-3-flash}                 & 0.50  & 3.00   & $\sim\computeCostCOT{0.50}{3.00}$  & $\sim\computeCostOS{0.50}{3.00}$ \\
                Gemini 2.5 Flash\cite{geminy-2.5-flash}             & 0.30  & 2.50   & $\sim\computeCostCOT{0.30}{2.50}$  & $\sim\computeCostOS{0.30}{2.50}$ \\
            \bottomrule
        \end{tabular}
    }
    \caption{Přibližné ceny vybraných cloudových API modelů. Náklady jsou vypočítány pro dva přístupy: chain-of-thought (CoT) a one-shot (OS), na základě průměrného počtu vstupních a výstupních tokenů pro jedno odevzdání.}
    \label{tab:cloud-api-costs}
\end{table}

Z pohledu nákladů se jako nejvýhodnější jeví modely GPT-5.4-nano a Claude 3.0 Haiku, které vychází na méně než \$0,002 za odevzdání při one-shot přístupu.
Použití chain-of-thought přístupu se třemi kroky analýzy je výrazně nákladnější: modely střední třídy (Claude 3.5 Haiku, GPT-5.4-mini) vychází na přibližně \$0,028--\$0,030 za odevzdání, oproti méně než \$0,004 při one-shot přístupu.
Výjimkou je Gemini 3.1 Flash-Lite, který i při chain-of-thought přístupu dosahuje nákladů okolo \$0,01 za odevzdání.

Nejvýkonnější modely (Claude 4.6 Opus, Claude 4.6 Sonnet, GPT-5.4) přesahují \$0,10 za odevzdání u CoT přístupu, přičemž Claude 4.6 Opus dosahuje nejvyšších nákladů \$0,174, což je pro rozsáhlé školní nasazení výrazně nákladnější.
Pro konkrétní volbu modelu bude rozhodující výsledek srovnávacího testování kvality výstupů.
Z pohledu nákladů vychází nejlépe menší modely, pokud jejich analytická kvalita pro dané požadavky je dostatečná.

\endinput
\subsection{On-premise řešení}
\label{subsec:deployment-on-premise}

On-premise nasazení znamená provozování jazykového modelu přímo na lokálních serverech, bez využití externích cloudových služeb.
V tomto případě je model hostován na serverové infrastruktuře a veškeré zpracování dat probíhá lokálně.

Výhodou on-premise řešení je především nezávislost na externích poskytovatelích služeb a absence poplatků za jednotlivé API požadavky.
Správci mají zároveň možnost plně kontrolovat konfiguraci modelu, způsob jeho aktualizace i způsob škálování výpočetních prostředků.
Na druhou stranu je nutné zajistit dostatečné hardwarové prostředky, zejména výkonné grafické akcelerátory (GPU), dostatek operační paměti a vhodnou infrastrukturu pro provoz a správu modelů.

Pro usnadnění provozu jazykových modelů v lokálním prostředí existuje několik softwarových nástrojů, které poskytují rozhraní pro jejich správu, spouštění a inferenci:

\begin{itemize}
    \item \textbf{Ollama}~\cite{ollama} -- nástroj zaměřený na jednoduché spouštění a správu jazykových modelů v lokálním prostředí.
    Poskytuje jednotné rozhraní pro práci s různými modely a umožňuje jejich snadné stažení, spuštění i aktualizaci.
    Součástí nástroje je také API, které umožňuje integraci modelů do externích aplikací.

    \item \textbf{Groq}~\cite{groq} -- platforma zaměřená na vysoce výkonnou inferenci jazykových modelů.
    Nabízí optimalizovanou infrastrukturu pro rychlé zpracování požadavků nad velkými modely.
    Přestože je Groq často využíván jako cloudová služba, jeho technologie je založena na specializovaném hardwaru navrženém pro efektivní provoz modelů.

    \item \textbf{vLLM}~\cite{vllm} -- knihovna určená pro efektivní provoz velkých jazykových modelů, která se zaměřuje zejména na optimalizaci paměťového využití a paralelního zpracování požadavků.
    Využívá techniky jako \textit{PagedAttention}, které umožňují efektivně sdílet paměť mezi více požadavky a tím výrazně zvyšují propustnost systému.

    \item \textbf{llama.cpp}~\cite{llamacpp} -- open-source implementace zaměřená na provoz modelů rodiny LLaMA a jejich variant i na běžných procesorech bez nutnosti využití GPU\@.
    Nástroj umožňuje provoz kvantizovaných verzí modelů, což výrazně snižuje paměťové nároky a umožňuje nasazení i na méně výkonném hardwaru.
\end{itemize}

Volba konkrétního nástroje závisí především na dostupném hardwaru, požadované rychlosti inferenčního procesu a způsobu integrace do existující infrastruktury systému Kelvin.
Zatímco nástroje jako \textit{vLLM} jsou vhodné pro provoz ve výkonném serverovém prostředí s GPU akcelerací, řešení jako \textit{llama.cpp} umožňují provoz modelů i na méně výkonném hardwaru.
Nástroj \textit{Ollama} pak poskytuje kompromis mezi jednoduchostí nasazení a možností integrace do aplikací prostřednictvím API\@.

\subsubsection*{Hardwarové nároky}

Pro on-premise nasazení je zcela zásadní dostupnost vhodného hardware.
Inference jazykových modelů je výpočetně náročná operace (viz \secref{subsec:llm-computational-cost}), proto je využití GPU s dostatečnou kapacitou paměti VRAM téměř nezbytné pro přijatelné výkony.
přibližné nároky na hardware pro různé modely lze vidět v \tabref{tab:memory-requirements}.

\TODO{Doplnit porovnání CPU vs GPU pro inferenci. Spustit ollama na argexu a změřit výkon pro různé modely.}

\subsubsection*{Limity}

On-premise nasazení přináší vlastní sadu omezení:

\begin{itemize}
    \item \textbf{Omezení dostupným hardware} -- kvalita a rychlost analýzy jsou přímou funkcí dostupného GPU.
    Na nedostatečném hardware je nutné volit menší nebo silněji kvantizované modely, což může snižovat kvalitu výstupů.

    \item \textbf{Škálovatelnost} -- cloudové API automaticky škáluje s nárůstem zátěže, kdežto lokální server má fixní kapacitu.
    Při hromadném odevzdávání se požadavky řadí do fronty a čekají na zpracování.

    \item \textbf{Správa modelu} -- nové verze modelů musí být ručně staženy a nasazeny.
    Aktualizace nebo výměna modelu vyžaduje přímou intervenci správce systému.
\end{itemize}

\subsubsection*{Ochrana dat}

Z pohledu ochrany dat je on-premise nasazení výrazně jednodušší situace.
Studentské zdrojové kódy a zadání úloh nikdy neopustí infrastrukturu školy.
Neexistuje třetí strana, které by musela být svěřena správa dat, a není potřeba uzavírat DPA s externím poskytovatelem.
Z hlediska GDPR jde o přímočarou variantu: správce dat (škola) je totožný s provozovatelem infrastruktury.

\subsubsection*{Náklady}

Náklady on-premise nasazení mají odlišnou strukturu než u cloudového API:

\begin{itemize}
    \item \textbf{Jednorázové pořizovací náklady} -- GPU server vhodný pro inferenci (např.\ NVIDIA RTX 4090 s 24 GB VRAM) se pohybuje v ceně od přibližně 60~000 Kč výše.
    Pro výkonnější varianty (NVIDIA A100) jsou náklady řádově vyšší (200~000--500~000 Kč).

    \item \textbf{Průběžné provozní náklady} -- elektřina (moderní GPU má spotřebu 150--400 W při zatížení), chlazení serverovny a lidská práce nutná pro správu a aktualizaci systému.

    \item \textbf{Nulové náklady za token} -- po zaplacení hardware jsou marginalní náklady na každé zpracované odevzdání prakticky nulové, bez ohledu na počet odevzdání.
\end{itemize}

Pro systém s malým počtem uživatelů a odevzdání může být on-premise pořizovací investice ekonomicky nevýhodná ve srovnání s cloudovým API\@.
Naopak pro systém se stovkami studentů zpracovávajícími tisíce odevzdání ročně se on-premise investice může vyplatit v horizontu 1--2 let.

\begin{table}[H]
    \centering
    \resizebox{\textwidth}{!}{%
        \begin{tabular}{lll}
            \toprule
            \textbf{Vlastnost}   & \textbf{Cloudové API}              & \textbf{On-premise nasazení} \\
            \midrule
            Hardwarové nároky    & Žádné (výpočet u poskytovatele)    & GPU server (10+ GB VRAM)     \\
            Pořizovací náklady   & Nulové                             & Vysoké (GPU hardware)        \\
            Průběžné náklady     & Za token (opakující se)            & Elektřina a správa           \\
            Škálovatelnost       & Automatická (rate limity)          & Fixní kapacita               \\
            Ochrana dat          & Data opouštějí školu               & Data zůstávají na škole      \\
            Kvalita modelů       & Nejnovější (GPT-4o, Claude Sonnet) & Závisí na hardware           \\
            Latence              & Závislá na internetu a zátěži API  & Závisí na hardware           \\
            Správa a aktualizace & Automatická (poskytovatel)         & Manuální (správce)           \\
            \bottomrule
        \end{tabular}
    }
    \caption{Srovnání cloudového API a on-premise nasazení jazykového modelu z hlediska klíčových vlastností}
    \label{tab:deployment-comparison}
\end{table}

\endinput

\endinput